\section{Nền tảng AIoT và xử lý tại biên}

Internet vạn vật (Internet of Things - IoT) là mô hình kết nối các đối tượng vật lý được nhúng cảm biến, bộ chấp hành và công nghệ truyền thông nhằm thu thập, trao đổi dữ liệu qua Internet \citeweb{ref_iotanalytics_2026}. Một hệ thống IoT truyền thống thường gồm ba lớp chính: lớp nhận thức thu thập dữ liệu từ môi trường vật lý, lớp mạng truyền dẫn dữ liệu, và lớp ứng dụng hoặc đám mây đảm nhiệm lưu trữ, phân tích và cung cấp dịch vụ cho người dùng. Trong bối cảnh hiện nay, quy mô kết nối toàn cầu tiếp tục tăng nhanh, với số thiết bị IoT được dự báo đạt khoảng 21,9 tỷ trong năm 2026 \citeweb{ref_iotinsider_2026}.

Sự mở rộng này làm bộc lộ hạn chế của mô hình IoT phụ thuộc hoàn toàn vào đám mây. Khi toàn bộ dữ liệu thô, đặc biệt là các tín hiệu có tần số lấy mẫu cao như rung động và âm thanh, đều được truyền lên máy chủ trung tâm, hệ thống phải đối mặt với áp lực băng thông, chi phí lưu trữ, độ trễ phản hồi và rủi ro riêng tư dữ liệu \cite{ref_1_4}. Các vấn đề này trở nên rõ rệt hơn trong bài toán giám sát tình trạng thiết bị, nơi hệ thống cần phản ứng nhanh trước các dấu hiệu bất thường về cơ khí.

Trí tuệ nhân tạo vạn vật (Artificial Intelligence of Things - AIoT) ra đời từ sự hội tụ giữa trí tuệ nhân tạo và IoT \citeweb{ref_1_2}. Thay vì chỉ thu thập và truyền dữ liệu thụ động, thiết bị AIoT có khả năng xử lý, phân tích và ra quyết định ngay tại hoặc gần nguồn phát sinh dữ liệu. Trong phạm vi đồ án, hướng tiếp cận này cho phép nút cảm biến gắn ngoài máy giặt tự trích xuất đặc trưng rung động và âm thanh, thực hiện suy luận bất thường tại chỗ, sau đó chỉ truyền kết quả cảnh báo hoặc dữ liệu đã rút gọn lên ứng dụng giám sát. Cách thiết kế này phù hợp với yêu cầu giảm tải đường truyền, giảm độ trễ và tăng tính độc lập của hệ thống khi giám sát các lỗi như mất cân bằng tải, lệch lồng giặt hoặc hư hỏng vòng bi \cite{ref_1_3}.

\subsection{Điện toán biên và suy luận tại thiết bị}

Điện toán biên (Edge Computing) là mô hình xử lý dữ liệu phân tán, trong đó dữ liệu được xử lý gần nguồn phát sinh thay vì phụ thuộc hoàn toàn vào hạ tầng đám mây \cite{ref_1_4}. So với xử lý tập trung, điện toán biên giúp giảm độ trễ truyền dẫn, tiết kiệm băng thông và hạn chế việc đưa dữ liệu nhạy cảm của người dùng ra khỏi thiết bị. Đây là cơ sở quan trọng cho các hệ thống giám sát thời gian thực trong môi trường gia dụng, nơi tín hiệu cảm biến có thể được sinh ra liên tục nhưng người dùng chỉ cần nhận các thông tin trạng thái hoặc cảnh báo đã được xử lý.

Kế thừa triết lý của điện toán biên, TinyML tập trung vào việc triển khai các mô hình học máy trên hệ thống nhúng và vi điều khiển có năng lượng thấp, bộ nhớ hạn chế và năng lực tính toán nhỏ hơn nhiều so với máy tính thông thường \cite{ref_1_1,ref_1_7}. Điểm khác biệt cốt lõi của TinyML nằm ở khả năng suy luận tại thiết bị. Mô hình sau huấn luyện được nén, lượng tử hóa và chuyển đổi sang định dạng phù hợp để thực thi trực tiếp trên vi điều khiển, qua đó giảm phụ thuộc vào kết nối Internet và máy chủ từ xa.

Trong đồ án này, vi điều khiển Seeed Studio XIAO ESP32-S3 kết hợp với TensorFlow Lite Micro được sử dụng làm nền tảng suy luận tại biên \citeweb{ref_tflite_micro}. Mô hình mạng nơ-ron tự mã hóa được lượng tử hóa về INT8 bằng kỹ thuật lượng tử hóa nhận thức huấn luyện, giúp giảm dung lượng lưu trữ và tăng hiệu quả tính toán trên phần cứng nhúng. Theo kết quả đo thực tế của firmware, tổng dung lượng ba mô hình cho ba pha vận hành là khoảng 95,3 KB, trong khi độ trễ suy luận trung bình đạt 6,4 ms. Các giá trị này cho thấy hệ thống có khả năng xử lý tại thiết bị với chi phí tài nguyên thấp, đồng thời vẫn duy trì được khả năng phản ứng thời gian thực.
